% GENERATED FILE. Edit canonical lessons, then run npm run generate:books.
% canonical-source-hash: fnv1a64:dca41f5fcdb659c7
% canonical-lessons: FR-C438-flaner, FR-C438-balader, FR-C438-deambuler, FR-C438-errer, FR-C438-arpenter

\chapter{Flâner en ville}
\label{ch:fr-flaner-en-ville}
\hlchaptermodality{438}

\begin{chapteropening}
\textbf{By the end of this chapter:} \emph{I can talk about wandering round a town: strolling, going for a walk, walking about, wandering and pacing the streets.}

Complete the last lesson of chapter 438: I can talk about wandering round a town: strolling, going for a walk, walking about, wandering and pacing the streets.
\end{chapteropening}

\section[flâner]{flâner — to stroll, to wander about}

\label{lesson:FR-C438-flaner}



\begin{quote}
Before the new one: say the French for the boulevard, then the French for the passer-by.
\end{quote}

\subsection*{You'll want to know: flâner}
\textbf{flâner} — \textquotedblleft{}to stroll, to wander about\textquotedblright{}.

\begin{quote}
Le dimanche, j'aime flâner au marché.
\end{quote}

\begin{quote}
\emph{On Sundays I like strolling round the market.}
\end{quote}

\emph{Un flâneur} — someone who strolls with no aim.

\begin{cousinweb}
\emph{flâner} came into French from a Norman dialect word, which is usually traced to Old Norse \emph{flana} 'to rush about'. EtymDB lists the verb without a source.
\end{cousinweb}

\begin{grammarlens}[title={What you've built}]
One more word for telling, explaining and describing at B1.
\end{grammarlens}

\subsection*{Your turn}
\begin{itemize}
\raggedright
  \item \emph{Say it:} \emph{flâner}
  \item \emph{Say it:} \emph{flâner}, once more
  \item \emph{Say it:} a sentence of your own with \emph{flâner}
  \item \emph{Recall:} say the French for the boulevard, then the French for the passer-by, then say \emph{flâner} again
\end{itemize}

\begin{tcolorbox}[breakable,colback=teal!4,colframe=teal!35!black,title={Before you move on}]
What does flâner mean? (\textquotedblleft{}To stroll, to wander about\textquotedblright{}.) Say it once more.
\end{tcolorbox}
\section[balader]{balader — to go for a walk (mostly reflexive: se balader); to take round}

\label{lesson:FR-C438-balader}



\begin{quote}
Before the new one: say the French for the passer-by, then the French for to stroll.
\end{quote}

\subsection*{You'll want to know: balader}
\textbf{balader} — \textquotedblleft{}to go for a walk (mostly reflexive: se balader); to take round\textquotedblright{}.

\begin{quote}
Après le dîner, nous nous baladons au bord de la rivière.
\end{quote}

\begin{quote}
\emph{After dinner we go for a walk by the river.}
\end{quote}

Mostly reflexive and informal: \emph{se balader}. \emph{Une balade} — a walk.

\begin{cousinweb}
\emph{balader} is built on \emph{balade} 'walk, outing', from \emph{ballade} 'dancing song', which goes back to Late Latin \emph{ballāre} 'to dance' (Lewis \& Short). EtymDB lists the verb without a source.
\end{cousinweb}

\begin{grammarlens}[title={What you've built}]
One more word for telling, explaining and describing at B1.
\end{grammarlens}

\subsection*{Your turn}
\begin{itemize}
\raggedright
  \item \emph{Say it:} \emph{balader}
  \item \emph{Say it:} \emph{balader}, once more
  \item \emph{Say it:} a sentence of your own with \emph{balader}
  \item \emph{Recall:} say the French for the passer-by, then the French for to stroll, then say \emph{balader} again
\end{itemize}

\begin{tcolorbox}[breakable,colback=teal!4,colframe=teal!35!black,title={Before you move on}]
What does balader mean? (\textquotedblleft{}To go for a walk (mostly reflexive: se balader); to take round\textquotedblright{}.) Say it once more.
\end{tcolorbox}
\section[déambuler]{déambuler — to walk about, to wander}

\label{lesson:FR-C438-deambuler}



\begin{quote}
Before the new one: say the French for to stroll, then the French for to go for a walk.
\end{quote}

\subsection*{You'll want to know: déambuler}
\textbf{déambuler} — \textquotedblleft{}to walk about, to wander\textquotedblright{}.

\begin{quote}
Les touristes déambulent entre les boutiques.
\end{quote}

\begin{quote}
\emph{The tourists wander about between the shops.}
\end{quote}

A bit literary.

\begin{cousinweb}
\emph{déambuler} is a learned borrowing from Latin \emph{deambulāre} 'to walk about', from \emph{dē-} and \emph{ambulāre} 'to walk' (EtymDB).
\end{cousinweb}

\begin{grammarlens}[title={What you've built}]
One more word for telling, explaining and describing at B1.
\end{grammarlens}

\subsection*{Your turn}
\begin{itemize}
\raggedright
  \item \emph{Say it:} \emph{déambuler}
  \item \emph{Say it:} \emph{déambuler}, once more
  \item \emph{Say it:} a sentence of your own with \emph{déambuler}
  \item \emph{Recall:} say the French for to stroll, then the French for to go for a walk, then say \emph{déambuler} again
\end{itemize}

\begin{tcolorbox}[breakable,colback=teal!4,colframe=teal!35!black,title={Before you move on}]
What does déambuler mean? (\textquotedblleft{}To walk about, to wander\textquotedblright{}.) Say it once more.
\end{tcolorbox}
\section[errer]{errer — to wander, to roam}

\label{lesson:FR-C438-errer}



\begin{quote}
Before the new one: say the French for to go for a walk, then the French for to walk about.
\end{quote}

\subsection*{You'll want to know: errer}
\textbf{errer} — \textquotedblleft{}to wander, to roam\textquotedblright{}.

\begin{quote}
Le chien errait seul dans la rue.
\end{quote}

\begin{quote}
\emph{The dog was wandering alone in the street.}
\end{quote}

Stronger than \emph{flâner}: no aim at all, sometimes lost.

\begin{cousinweb}
\emph{errer} is inherited through Old French \emph{errer} from Latin \emph{errāre} 'to wander, to go astray' (EtymDB).
\end{cousinweb}

\begin{grammarlens}[title={What you've built}]
One more word for telling, explaining and describing at B1.
\end{grammarlens}

\subsection*{Your turn}
\begin{itemize}
\raggedright
  \item \emph{Say it:} \emph{errer}
  \item \emph{Say it:} \emph{errer}, once more
  \item \emph{Say it:} a sentence of your own with \emph{errer}
  \item \emph{Recall:} say the French for to go for a walk, then the French for to walk about, then say \emph{errer} again
\end{itemize}

\begin{tcolorbox}[breakable,colback=teal!4,colframe=teal!35!black,title={Before you move on}]
What does errer mean? (\textquotedblleft{}To wander, to roam\textquotedblright{}.) Say it once more.
\end{tcolorbox}
\section[arpenter]{arpenter — to pace up and down; to stride along}

\label{lesson:FR-C438-arpenter}



\begin{quote}
Before the new one: say the French for to walk about, then the French for to wander.
\end{quote}

\subsection*{You'll want to know: arpenter}
\textbf{arpenter} — \textquotedblleft{}to pace up and down; to stride along\textquotedblright{}.

\begin{quote}
Pendant deux jours, nous avons arpenté la ville à pied.
\end{quote}

\begin{quote}
\emph{For two days we walked the town from end to end.}
\end{quote}

Also \emph{arpenter le couloir} — to pace up and down the corridor.

\begin{cousinweb}
\emph{arpenter} is built on \emph{arpent}, an old measure of land (EtymDB), from Latin \emph{arepennis} 'half an acre', a word Latin took from Gaulish (Lewis \& Short): first to measure a field by walking it.
\end{cousinweb}

\begin{grammarlens}[title={What you've built}]
One more word for telling, explaining and describing at B1.
\end{grammarlens}

\subsection*{Your turn}
\begin{itemize}
\raggedright
  \item \emph{Say it:} \emph{arpenter}
  \item \emph{Say it:} \emph{arpenter}, once more
  \item \emph{Say it:} a sentence of your own with \emph{arpenter}
  \item \emph{Recall:} say the French for to walk about, then the French for to wander, then say \emph{arpenter} again
\end{itemize}

\begin{tcolorbox}[breakable,colback=teal!4,colframe=teal!35!black,title={Before you move on}]
What does arpenter mean? (\textquotedblleft{}To pace up and down; to stride along\textquotedblright{}.) Say it once more.
\end{tcolorbox}
